\subsection{Hubbard--Stratonovich transformation for \texorpdfstring{$N$}{N} scalar fields}
\label{p:app:scalar-hs}


In this appendix, we perform the Hubbard--Stratonovich transformation for a theory of several scalar fields with a quartic interaction.
We see on this example how the auxiliary fields are organized when there are several physical fields: the auxiliary field carries a pair of indices, as $\Sigma_2$ carries two strings.
This is the model of the paper's introduction with constant couplings, which underlies the analogy used in \cref{p:sec:hs:quartic}.


\subparagraph{Action of the scalar fields.}

Consider $N$ scalar fields $\phi_i$, $i = 1, \ldots, N$, with the action:
\begin{equation}
\label{p:eq:shs:action}
S
  = \int \dd^d x \left[
    - \frac{1}{2} \, \partial_\mu \phi_i \, \partial^\mu \phi_i
    - \frac{1}{2} \, M_{ij} \phi_i \phi_j
    + \frac{1}{3!} \, W_{ijk} \phi_i \phi_j \phi_k
    + \frac{1}{4!} \, V_{ijk\ell} \phi_i \phi_j \phi_k \phi_\ell
  \right],
\end{equation}
where repeated indices are summed over.
The mass matrix $M_{ij}$ and the coupling tensors $W_{ijk}$ and $V_{ijk\ell}$ are constant and completely symmetric.
We write $S_4$ for the quartic term.


\subparagraph{Auxiliary fields.}

The quartic term is the square of the products $\phi_i \phi_j$ with respect to the tensor $V_{ijk\ell}$.
We introduce one auxiliary field for each such product, that is $N (N + 1) / 2$ fields $\sigma_{ij} = \sigma_{ji}$, and add to the action the term~\cite{Stratonovich:1957:MethodCalculatingQuantum, Hubbard:1959:CalculationPartitionFunctions, Weinberg:2005:QuantumTheoryFields-2}:
\begin{equation}
\label{p:eq:shs:DeltaS4}
\Delta S_4
  := - \frac{1}{4!} \int \dd^d x \, V_{ijk\ell} \, \big(
    \sigma_{ij} - \phi_i \phi_j
  \big) \big(
    \sigma_{k\ell} - \phi_k \phi_\ell
  \big),
\end{equation}
which vanishes when $\sigma_{ij} = \phi_i \phi_j$.
We expand the square and use $V_{ijk\ell} = V_{k\ell ij}$ for the cross term, so that the quartic term of the action becomes:
\begin{equation}
\label{p:eq:shs:S4tilde}
\caction[4]
  := S_4 + \Delta S_4
  = \int \dd^d x \left[
    - \frac{1}{4!} \, V_{ijk\ell} \, \sigma_{ij} \sigma_{k\ell}
    + \frac{1}{12} \, V_{ijk\ell} \, \sigma_{ij} \phi_k \phi_\ell
  \right].
\end{equation}
The term quartic in $\phi_i$ has disappeared.
The kinetic, mass and cubic terms are untouched, so that the action
\begin{equation}
\label{p:eq:shs:action-tilde}
\begin{aligned}
  \caction
    := S + \Delta S_4
    = \int \dd^d x \bigg[
      &
      - \frac{1}{2} \, \partial_\mu \phi_i \, \partial^\mu \phi_i
      - \frac{1}{2} \, M_{ij} \phi_i \phi_j
      + \frac{1}{3!} \, W_{ijk} \phi_i \phi_j \phi_k
      \\ &
      + \frac{1}{12} \, V_{ijk\ell} \, \sigma_{ij} \phi_k \phi_\ell
      - \frac{1}{4!} \, V_{ijk\ell} \, \sigma_{ij} \sigma_{k\ell}
    \bigg]
\end{aligned}
\end{equation}
is cubic in $(\phi_i, \sigma_{ij})$.
For a single field, $N = 1$, the integrand of the quartic term~\eqref{p:eq:shs:S4tilde} reads:
\begin{equation}
- \frac{1}{4!} \, V \sigma^2
+ \frac{1}{12} \, V \sigma \phi^2.
\end{equation}
This is the transformation of the introduction at constant coupling, with the normalization $1 / 4!$ of the vertices.

The fields $\sigma_{ij}$ have no kinetic term, and their equation of motion is algebraic:
\begin{equation}
\label{p:eq:shs:eom-sigma}
V_{ijk\ell} \, \big( \sigma_{k\ell} - \phi_k \phi_\ell \big)
  = 0.
\end{equation}
When $V_{ijk\ell}$ is invertible as a matrix acting on the symmetric pairs of indices $(ij)$ and $(k\ell)$ (this is the matrix $V_{IJ}$ defined below), the solution is:
\begin{equation}
\label{p:eq:shs:solution}
\sigma_{ij}
  = \phi_i \phi_j,
\end{equation}
on which $\Delta S_4$ vanishes and $\caction$ reduces to the original action $S$.


\subparagraph{Pair-index basis.}

The pair of indices of $\sigma_{ij}$ can be traded for a single index.
Let $\lambda_{I;ij}$, $I = 1, \ldots, N (N + 1) / 2$, be an orthonormal and complete basis of the real symmetric $N \times N$ matrices:
\begin{equation}
\label{p:eq:shs:basis}
\lambda_{I;ij} \lambda_{J;ij}
  = \delta_{IJ},
\qquad
\lambda_{I;ij} \lambda_{I;k\ell}
  = \frac{1}{2} \, \big( \delta_{ik} \delta_{j\ell} + \delta_{i\ell} \delta_{jk} \big).
\end{equation}
The right-hand side of the completeness relation is the projector on symmetric matrices.
We define:
\begin{equation}
\label{p:eq:shs:pair-index}
\sigma_I
  := \lambda_{I;ij} \, \sigma_{ij},
\qquad
V_{IJ}
  := \lambda_{I;ij} \lambda_{J;k\ell} \, V_{ijk\ell},
\qquad
V_{I;k\ell}
  := \lambda_{I;ij} \, V_{ijk\ell}.
\end{equation}
Since $\sigma_{ij}$ and $V_{ijk\ell}$ are symmetric in each pair of indices, the completeness relation gives $\sigma_{ij} = \lambda_{I;ij} \sigma_I$ and $V_{I;k\ell} = V_{IJ} \lambda_{J;k\ell}$.
The quartic term~\eqref{p:eq:shs:S4tilde} then becomes:
\begin{equation}
\label{p:eq:shs:S4tilde-pair}
\caction[4]
  = \int \dd^d x \left[
    - \frac{1}{4!} \, V_{IJ} \, \sigma_I \sigma_J
    + \frac{1}{12} \, V_{I;k\ell} \, \sigma_I \phi_k \phi_\ell
  \right].
\end{equation}

The matrix $V_{IJ}$ is real and symmetric, so we can diagonalize it by an orthogonal change of the basis $\lambda_{I;ij}$, which puts the quadratic term of the $\sigma_I$ in diagonal form.
We assume that all its eigenvalues are nonzero.
Otherwise, the auxiliary fields $\sigma_I$ along the eigenvectors with zero eigenvalue appear neither in the quadratic term nor in the coupling to $\phi_i$, since $V_{I;k\ell} = V_{IJ} \lambda_{J;k\ell}$.
They do not enter the action, so we leave them out and perform the transformation on the subset of the $\sigma_I$ that does enter.


\subparagraph{Four-point amplitude.}

We read the Feynman rules off $\caction$ as in \cref{p:sec:hs:quartic}: a vertex factor is the derivative of the action with respect to the fields on its legs, and a propagator is minus the inverse of the Hessian.
The auxiliary sector~\eqref{p:eq:shs:S4tilde-pair} contributes the propagator of $\sigma_I$ and the vertex coupling $\sigma_I$ to $\phi_k \phi_\ell$, respectively:
\begin{equation}
\label{p:eq:shs:feynman-rules}
12 \, (V^{-1})_{IJ},
\qquad
\frac{1}{3!} \, V_{I;k\ell}.
\end{equation}

The exchange of $\sigma$ between two such vertices, with the external fields $\phi_i, \phi_j$ on one vertex and $\phi_k, \phi_\ell$ on the other, gives:
\begin{equation}
\label{p:eq:shs:channel}
\frac{1}{3!} \, V_{I;ij} \;
12 \, (V^{-1})_{IJ} \;
\frac{1}{3!} \, V_{J;k\ell}
  = \frac{1}{3} \, \lambda_{K;ij} V_{KI} (V^{-1})_{IJ} V_{JL} \lambda_{L;k\ell}
  = \frac{1}{3} \, V_{ijk\ell},
\end{equation}
where the first equality uses $V_{I;ij} = V_{IK} \lambda_{K;ij}$ and the second $V_{ijk\ell} = \lambda_{K;ij} V_{KL} \lambda_{L;k\ell}$, which follows from the completeness relation.
We sum over the three ways of distributing the four external fields on the two vertices and use the complete symmetry of $V_{ijk\ell}$.
The four-point amplitude is then:
\begin{equation}
\label{p:eq:shs:A4}
A_{ijk\ell}
  = \frac{1}{3} \, \big( V_{ijk\ell} + V_{ikj\ell} + V_{i\ell jk} \big) + \cdots
  = V_{ijk\ell} + \cdots,
\end{equation}
where the dots are the $s$-, $t$- and $u$-channel exchanges of $\phi$ built from the cubic coupling $W_{ijk}$, which the transformation leaves untouched.
The exchange of $\sigma$ thus reproduces the contact vertex $V_{ijk\ell}$ of the original action~\eqref{p:eq:shs:action}.


\subparagraph{Relation to the string field.}

We obtain the construction of \cref{p:sec:hs:quartic} from this one by replacing the products $\phi_i \phi_j$ by the $2$-string state $\Psi^2$, the tensor $V_{ijk\ell}$ by $\gs^2 \vertex{4}$ (so that $V_{IJ}$ becomes $\gs^2 \sprod{2}{2}$), the auxiliary field $\sigma_{ij}$ by $\Sigma_2 \in \Hn{2}$, and the propagator $12 \, (V^{-1})_{IJ}$ by $12 \gs^{-2} \, \sprod{2}{2}^{-1}$.



\subparagraph{Degenerate coupling matrix.}

To couple the directions of zero eigenvalue of $V_{IJ}$, left out above, we now add a quintic term to the action~\eqref{p:eq:shs:action}, with a constant and completely symmetric coupling $U_{ijk\ell m}$:
\begin{equation}
\label{p:eq:shs:S5}
S_5
  = \frac{1}{5!} \int \dd^d x \, U_{ijk\ell m} \, \phi_i \phi_j \phi_k \phi_\ell \phi_m.
\end{equation}
As in \cref{p:sec:hs:quintic}, we replace two pairs of fields by the auxiliary field, which turns the quintic term into a cubic coupling:
\begin{equation}
\label{p:eq:shs:S5tilde}
\caction[5]
  := \frac{1}{5!} \int \dd^d x \, U_{IJm} \, \sigma_I \sigma_J \phi_m,
\qquad
U_{IJm}
  := \lambda_{I;ij} \lambda_{J;k\ell} \, U_{ijk\ell m},
\end{equation}
which agrees with $S_5$ at $\sigma_{ij} = \phi_i \phi_j$.
When the matrix $V_{IJ}$ is invertible, eliminating $\sigma_I$ from $\caction + \caction[5]$ gives back $S + S_5$ up to terms of second order in $U_{ijk\ell m}$, the exchange of $\sigma$ between two quintic couplings, whose string counterpart is cancelled at the sextic order (\cref{p:sec:hs:sextic}).

We now suppose that $V_{IJ}$ has null vectors, $V_{IJ} \xi^J = 0$, and we work in the basis $\lambda_{I;ij}$ that diagonalizes it.
The index $I$ then splits into the directions of nonzero eigenvalue, which span the orthogonal complement $C$ of the null space, and the null directions, which span the null space $K$.
We write $\pi_C$ and $\pi_K$ for the projections on $C$ along $K$ and on $K$ along $C$.
The null components of $\sigma_I$ drop out of the quartic terms~\eqref{p:eq:shs:S4tilde-pair}, as noted above, but they enter the quintic coupling~\eqref{p:eq:shs:S5tilde}.
Under the shift $\sigma_I \to \sigma_I + \xi^I$, $\caction[5]$ changes by:
\begin{equation}
\label{p:eq:shs:null-shift}
\frac{1}{5!} \int \dd^d x \, \big(
  2 \, U_{IJm} \, \xi^I \sigma_J \phi_m
  + U_{IJm} \, \xi^I \xi^J \phi_m
\big).
\end{equation}
For this to vanish for all $\sigma_I$, the term linear in $\sigma_J$ must vanish, which then implies that the quadratic one vanishes.
The null components therefore decouple if and only if the quintic coupling annihilates the null vectors, and they can then be left out as before:
\begin{equation}
\label{p:eq:shs:kernel-condition}
\forall \xi \in K, \;
\forall k, \ell, m:
\quad
\xi^I \, U_{I;k\ell m}
  = 0,
\qquad
U_{I;k\ell m}
  := \lambda_{I;ij} \, U_{ijk\ell m}.
\end{equation}
This is the condition~\eqref{p:eq:hs5:kernel-condition} of the string construction, with $\xi$ for $N_2$ and $U_{ijk\ell m}$ for $\vertex{5}$.

When this condition fails, the null components of $\sigma_I$ have no quadratic term in $\caction$ but they have a source, the term of~\eqref{p:eq:shs:S5tilde} linear in them.
  We could keep the null components as fields, but their quadratic form in $\caction[5]$ is proportional to $\phi_m$, so that around $\phi_i = 0$ they have neither a propagator nor a perturbative expansion.
  If we leave them out instead, that is, if we take $\sigma_I$ in $C$, we get an action of the form~\eqref{p:eq:shs:action-tilde} with $V_{IJ}$ restricted to $C$.
  Since $V_{IJ}$ is invertible on $C$, $\sigma_I$ can be eliminated as before.
  The eliminated action reproduces the quartic term.
  However, at first order in $U_{ijk\ell m}$, it contains $U_{IJm} \, (\pi_C P)^I (\pi_C P)^J \phi_m$ in place of the quintic term $U_{IJm} \, P_I P_J \phi_m$, where $P_I := \lambda_{I;ij} \, \phi_i \phi_j$.
  The missing terms are those with a null direction in at least one of the two pair indices; they depend on the complement and cannot vanish for all complements.
  Indeed, we could choose any other complement of $K$ in place of the span of the eigenvectors, and we would get a different eliminated action.

We can repair the construction with a multiplier on the null directions, as in the hybrid form of \cref{p:sec:hs:sextic}.
  We keep all the components of $\sigma_I$, and we add to the action a field $\theta_I$ on the null directions, paired with the null component of $\sigma_I - P_I$:
  \begin{equation}
  \label{p:eq:shs:multiplier}
  \caction + \caction[5]
    + \int \dd^d x \, \theta_I \, \big( \pi_K (\sigma - P) \big)^I.
\end{equation}
  Since $\theta_I P_I$ is cubic, every term has at most three fields and the action stays cubic.
  The equation of motion of $\theta_I$ sets the null components of $\sigma_I$ to those of $P_I$.
  The equation of motion of the null components fixes $\theta_I$, which appears nowhere else.
  Finally, the equation of motion of the components along $C$ is the one of the action with $\sigma_I$ in $C$, with a source shifted by the quintic coupling evaluated on the null component of $P_I$.
  Eliminating all the fields gives back $S + S_5$ up to terms of second order in $U_{ijk\ell m}$, for every choice of complement.
  Indeed, the null components supply the part of $P_I$ that the complement was missing.
  The multiplier vanishes on the solution if and only if the quintic coupling annihilates the null vectors, in which case it is not needed.
  We see on this degenerate case how the multiplier replaces the quotient when a vertex does not vanish on the kernel of the quadratic form, in a setting where the linear algebra is explicit.
  For the string vertices, a multiplier on the null directions is the fallback for an accidental null vector of $\sprod{2}{2}$, though we have no indication that one exists (\cref{p:sec:hs:quintic}).
